% GENERATED FILE. Edit canonical lessons, then run npm run generate:books.
% canonical-source-hash: fnv1a64:a5abb248a5f28bdc
% canonical-lessons: GU-C155-uthadvu, GU-C155-taiyar-thavu, GU-C155-olvu, GU-C155-attar, GU-C155-lagadvu

\chapter{Getting ready in the morning}
\label{ch:gu-getting-ready-in-the-morning}
\hlchaptermodality{155}

\begin{chapteropening}
\textbf{By the end of this chapter:} \emph{I can say how I get ready in the morning.}

Complete the last lesson of chapter 155: I can say how I get ready in the morning.
\end{chapteropening}

\section[\texorpdfstring{\gu{ઉઠાડવું} (uṭhāḍvũ)}{ઉઠાડવું (uṭhāḍvũ)}]{\gu{ઉઠાડવું} (uṭhāḍvũ) — to wake someone, to get someone up}

\label{lesson:GU-C155-uthadvu}



\begin{quote}
Before the new one: say the Gujarati for to take off, then the Gujarati for to wrap oneself in.
\end{quote}

\subsection*{You'll want to know: \gu{ઉઠાડવું}}
\textbf{\gu{ઉઠાડવું}} — \emph{uṭhāḍvũ} — \textquotedblleft{}to wake someone, to get someone up\textquotedblright{}.

\begin{quote}
\textbf{\gu{માતા} \gu{મને} \gu{સવારે} \gu{ઉઠાડે} \gu{છે}.} — \emph{mātā mane savāre uṭhāḍe che} — \textquotedblleft{}Mother wakes me in the morning.\textquotedblright{}
\end{quote}

Built on \textbf{\gu{ઉઠવું}}, \emph{uṭhvũ}, to get up: this verb is to get someone else up.

\begin{cousinweb}
\textbf{\gu{ઉઠાડવું}} (\emph{uṭhāḍvũ}) is built on the inherited \emph{uṭhvũ}, \textquotedblleft{}to get up\textquotedblright{}, with the causative \emph{-āḍ-}. CDIAL traces \emph{uṭhvũ} to \emph{*utsthāti}, \textquotedblleft{}stands up\textquotedblright{}, a reconstructed form (CDIAL 1900); it lists no Gujarati form of the causative itself.
\end{cousinweb}

\begin{grammarlens}[title={What you've built}]
One more word for everyday life at A2.
\end{grammarlens}

\subsection*{Your turn}
\begin{itemize}
\raggedright
  \item \emph{Say it:} \emph{uṭhāḍvũ}
  \item \emph{Say it:} \emph{uṭhāḍvũ}, once more
  \item \emph{Say it:} \emph{mātā mane savāre uṭhāḍe che} — Mother wakes me in the morning.
  \item \emph{Recall:} say the Gujarati for to take off, then the Gujarati for to wrap oneself in, then say \emph{uṭhāḍvũ} again
\end{itemize}

\begin{tcolorbox}[breakable,colback=teal!4,colframe=teal!35!black,title={Before you move on}]
What does \gu{ઉઠાડવું} mean? (\textquotedblleft{}To wake someone, to get someone up\textquotedblright{}.) Say it once more.
\end{tcolorbox}
\section[\texorpdfstring{\gu{તૈયાર} \gu{થવું} (taiyār thavũ)}{તૈયાર થવું (taiyār thavũ)}]{\gu{તૈયાર} \gu{થવું} (taiyār thavũ) — to get ready}

\label{lesson:GU-C155-taiyar-thavu}



\begin{quote}
Before the new one: say the Gujarati for to wrap oneself in, then the Gujarati for to wake someone.
\end{quote}

\subsection*{You'll want to know: \gu{તૈયાર} \gu{થવું}}
\textbf{\gu{તૈયાર} \gu{થવું}} — \emph{taiyār thavũ} — \textquotedblleft{}to get ready\textquotedblright{}.

\begin{quote}
\textbf{\gu{હું} \gu{સવારે} \gu{જલદી} \gu{તૈયાર} \gu{થાઉં} \gu{છું}.} — \emph{hũ savāre jaldī taiyār thāũ chũ} — \textquotedblleft{}I get ready quickly in the morning.\textquotedblright{}
\end{quote}

\begin{cousinweb}
\textbf{\gu{તૈયાર} \gu{થવું}} (\emph{taiyār thavũ}) is built from \emph{taiyār}, \textquotedblleft{}ready\textquotedblright{}, a loan from Persian \emph{tayyār} (EtymDB lists the Persian word), and the inherited verb \emph{thavũ}, \textquotedblleft{}to become\textquotedblright{}.
\end{cousinweb}

\begin{grammarlens}[title={What you've built}]
One more word for everyday life at A2.
\end{grammarlens}

\subsection*{Your turn}
\begin{itemize}
\raggedright
  \item \emph{Say it:} \emph{taiyār thavũ}
  \item \emph{Say it:} \emph{taiyār thavũ}, once more
  \item \emph{Say it:} \emph{hũ savāre jaldī taiyār thāũ chũ} — I get ready quickly in the morning.
  \item \emph{Recall:} say the Gujarati for to wrap oneself in, then the Gujarati for to wake someone, then say \emph{taiyār thavũ} again
\end{itemize}

\begin{tcolorbox}[breakable,colback=teal!4,colframe=teal!35!black,title={Before you move on}]
What does \gu{તૈયાર} \gu{થવું} mean? (\textquotedblleft{}To get ready\textquotedblright{}.) Say it once more.
\end{tcolorbox}
\section[\texorpdfstring{\gu{ઓળવું} (oḷvũ)}{ઓળવું (oḷvũ)}]{\gu{ઓળવું} (oḷvũ) — to comb one's hair, to comb}

\label{lesson:GU-C155-olvu}



\begin{quote}
Before the new one: say the Gujarati for to wake someone, then the Gujarati for to get ready.
\end{quote}

\subsection*{You'll want to know: \gu{ઓળવું}}
\textbf{\gu{ઓળવું}} — \emph{oḷvũ} — \textquotedblleft{}to comb one's hair, to comb\textquotedblright{}.

\begin{quote}
\textbf{\gu{તે} \gu{રોજ} \gu{વાળ} \gu{ઓળે} \gu{છે}.} — \emph{te roj vāḷ oḷe che} — \textquotedblleft{}She combs her hair every day.\textquotedblright{}
\end{quote}

\begin{cousinweb}
The origin of \textbf{\gu{ઓળવું}} (\emph{oḷvũ}) is not settled. CDIAL lists no Gujarati verb of this form with the sense of combing, and EtymDB has no entry for it.
\end{cousinweb}

\begin{grammarlens}[title={What you've built}]
One more word for everyday life at A2.
\end{grammarlens}

\subsection*{Your turn}
\begin{itemize}
\raggedright
  \item \emph{Say it:} \emph{oḷvũ}
  \item \emph{Say it:} \emph{oḷvũ}, once more
  \item \emph{Say it:} \emph{te roj vāḷ oḷe che} — She combs her hair every day.
  \item \emph{Recall:} say the Gujarati for to wake someone, then the Gujarati for to get ready, then say \emph{oḷvũ} again
\end{itemize}

\begin{tcolorbox}[breakable,colback=teal!4,colframe=teal!35!black,title={Before you move on}]
What does \gu{ઓળવું} mean? (\textquotedblleft{}To comb one's hair, to comb\textquotedblright{}.) Say it once more.
\end{tcolorbox}
\section[\texorpdfstring{\gu{અત્તર} (attar)}{અત્તર (attar)}]{\gu{અત્તર} (attar) — perfume, attar}

\label{lesson:GU-C155-attar}



\begin{quote}
Before the new one: say the Gujarati for to get ready, then the Gujarati for to comb one's hair.
\end{quote}

\subsection*{You'll want to know: \gu{અત્તર}}
\textbf{\gu{અત્તર}} — \emph{attar} — \textquotedblleft{}perfume, attar\textquotedblright{}.

\begin{quote}
\textbf{\gu{આ} \gu{અત્તર} \gu{મોંઘું} \gu{છે}.} — \emph{ā attar moṅghũ che} — \textquotedblleft{}This perfume is expensive.\textquotedblright{}
\end{quote}

\begin{cousinweb}
\textbf{\gu{અત્તર}} (\emph{attar}) is a loan from Arabic \emph{ʿiṭr}, \textquotedblleft{}perfume\textquotedblright{}, which reached Gujarati through Persian; EtymDB lists the Persian word, glossed \textquotedblleft{}scent\textquotedblright{}. Its Arabic root is \emph{ʿ-ṭ-r} (Lane).
\end{cousinweb}

\begin{grammarlens}[title={What you've built}]
One more word for everyday life at A2.
\end{grammarlens}

\subsection*{Your turn}
\begin{itemize}
\raggedright
  \item \emph{Say it:} \emph{attar}
  \item \emph{Say it:} \emph{attar}, once more
  \item \emph{Say it:} \emph{ā attar moṅghũ che} — This perfume is expensive.
  \item \emph{Recall:} say the Gujarati for to get ready, then the Gujarati for to comb one's hair, then say \emph{attar} again
\end{itemize}

\begin{tcolorbox}[breakable,colback=teal!4,colframe=teal!35!black,title={Before you move on}]
What does \gu{અત્તર} mean? (\textquotedblleft{}Perfume, attar\textquotedblright{}.) Say it once more.
\end{tcolorbox}
\section[\texorpdfstring{\gu{લગાડવું} (lagāḍvũ)}{લગાડવું (lagāḍvũ)}]{\gu{લગાડવું} (lagāḍvũ) — to apply, to put on}

\label{lesson:GU-C155-lagadvu}



\begin{quote}
Before the new one: say the Gujarati for to comb one's hair, then the Gujarati for perfume.
\end{quote}

\subsection*{You'll want to know: \gu{લગાડવું}}
\textbf{\gu{લગાડવું}} — \emph{lagāḍvũ} — \textquotedblleft{}to apply, to put on\textquotedblright{}.

\begin{quote}
\textbf{\gu{તે} \gu{વાળમાં} \gu{તેલ} \gu{લગાડે} \gu{છે}.} — \emph{te vāḷmā̃ tel lagāḍe che} — \textquotedblleft{}She puts oil in her hair.\textquotedblright{}
\end{quote}

\begin{cousinweb}
\textbf{\gu{લગાડવું}} (\emph{lagāḍvũ}) is built on the inherited \emph{lāgvũ}, \textquotedblleft{}to touch, to come into contact\textquotedblright{}, with the causative \emph{-āḍ-}. CDIAL traces \emph{lāgvũ} to Sanskrit \emph{lagyati}, \textquotedblleft{}adheres to\textquotedblright{} (CDIAL 10895).
\end{cousinweb}

\begin{grammarlens}[title={What you've built}]
One more word for everyday life at A2.
\end{grammarlens}

\subsection*{Your turn}
\begin{itemize}
\raggedright
  \item \emph{Say it:} \emph{lagāḍvũ}
  \item \emph{Say it:} \emph{lagāḍvũ}, once more
  \item \emph{Say it:} \emph{te vāḷmā̃ tel lagāḍe che} — She puts oil in her hair.
  \item \emph{Recall:} say the Gujarati for to comb one's hair, then the Gujarati for perfume, then say \emph{lagāḍvũ} again
\end{itemize}

\begin{tcolorbox}[breakable,colback=teal!4,colframe=teal!35!black,title={Before you move on}]
What does \gu{લગાડવું} mean? (\textquotedblleft{}To apply, to put on\textquotedblright{}.) Say it once more.
\end{tcolorbox}
